\subsection{Discussion}
\begin{itemize}
  \item \textbf{All adaptive controllers are far better than the fixed plan.} On the unseen days
        the fixed plan gives a mean wait of 485\,s; the actuated controller, the camera rule and
        the camera-view DQN give 98\,s, 63\,s and 104\,s. They also move about 11\% more vehicles
        per hour and use about 35\% less fuel per vehicle.
  \item \textbf{The camera rule is the best controller.} With only a 75\,m camera view it waits
        35\% less than the loop-sensor actuated controller on the unseen days (63\,s against
        98\,s), 25\% less on the unseen junction (40\,s against 53\,s) and 43\% less with free
        left turns. At 1.3$\times$ demand it equals the actuated controller (507\,s against 504\,s)
        and has a higher throughput. Its advantage comes from two choices: it skips roads with no
        queue, and it ranks roads by the PCU of the queue and the longest wait.
  \item \textbf{The camera-view DQN is close to the actuated controller.} On the unseen junction it
        waits 49\,s against 53\,s, and on the unseen days 104\,s against 98\,s. The result changes
        between training seeds (standard deviation 28\,s on the unseen days), so the mean of
        several seeds is reported.
  \item \textbf{The DQN does not handle a change of road layout.} With free left turns, which it
        never saw in training, its wait rises to 239\,s, while the two rules do not need training
        and stay good. A learned controller must train on the layout where it will run.
  \item \textbf{Buses and trucks are not left behind.} For every adaptive controller, buses wait
        less than cars and trucks wait about as long as cars (Table~\ref{tab:class}).
\end{itemize}
